\documentclass[10pt,a4paper]{article}

% --- XeLaTeX: fonts + Russian ---
\usepackage{fontspec}
\usepackage{polyglossia}
\setdefaultlanguage{russian}
\setotherlanguage{english}
\setmainfont{PT Sans}
\setmonofont[Scale=MatchLowercase]{PT Mono}
\newfontfamily\cyrillicfonttt[Scale=MatchLowercase]{PT Mono}

% --- layout ---
\usepackage[a4paper,top=0.75cm,bottom=0.75cm,left=1.3cm,right=1.3cm]{geometry}
\usepackage{titlesec}
\usepackage{enumitem}
\usepackage{xcolor}
\usepackage[hidelinks]{hyperref}

\definecolor{rule}{HTML}{333333}

\titleformat{\section}{\normalsize\bfseries}{}{0em}{\MakeUppercase}[{\color{rule}\titlerule[0.6pt]}]
\titlespacing*{\section}{0pt}{6pt}{3pt}

\setlist[itemize]{leftmargin=1.15em, labelsep=0.5em, itemsep=1.6pt, topsep=2pt, parsep=0pt, before=\vspace{1pt}}
\renewcommand{\labelitemi}{\textbullet}

\setlength{\parindent}{0pt}
\setlength{\parskip}{0pt}
\pagestyle{empty}
\linespread{1.0}
\tolerance=1500
\emergencystretch=2.5em

\newcommand{\role}[2]{\textbf{#1}\hfill #2\par}

\begin{document}

% ===== Header =====
\begin{center}
  {\LARGE\bfseries Михайловский Илья Евгеньевич}\\[2pt]
  {\large ML Engineer (NLP / LLM)}\\[3pt]
  Москва \textperiodcentered{} офис / гибрид / удалённо \textperiodcentered{} английский C1\\[2pt]
  Telegram: \href{https://t.me/Iluuuaaa}{@Iluuuaaa} \textperiodcentered{}
  \href{mailto:ilyamihailovsy@gmail.com}{ilyamihailovsy@gmail.com} \textperiodcentered{}
  \href{https://github.com/iluuua}{github.com/iluuua}
\end{center}

% ===== Summary =====
\section{О себе}
Студент ПМИ (Московский политех, выпуск 2029). С 2024 года пишу Python в продакшене, с начала 2026-го веду
LLM-слой B2B-продукта: конвейер из 8 вызовов модели на входящее сообщение, strict JSON-schema output, агент
с реестром инструментов и MCP, роутинг моделей с фолбэками, трассировка стоимости вызовов. Модели пишу с нуля
на PyTorch: финалист AIDAO 2025 (ФКН НИУ ВШЭ и Яндекс Образование), occupancy в bird's-eye view по камерам
беспилотника --- IoU 0.5606 при 0.564 у победителя. Алгоритмическая база --- C++, олимпиадное программирование.

% ===== Skills =====
\section{Навыки}
\begin{itemize}
  \item \textbf{ML / DL:} PyTorch, NumPy, pandas, matplotlib; обучение с нуля и transfer learning, mixed precision,
        аугментации, подбор лосса под метрику, калибровка порога, дисбаланс классов; ResNet / FPN, сегментация, IoU
  \item \textbf{NLP / LLM:} LLM-пайплайны в проде, prompt engineering, strict JSON-schema output, классификация и
        extraction, RAG, эмбеддинги + pgvector, LLM-агенты и tool calling, MCP, роутинг моделей и fallback,
        оценка cost / latency
  \item \textbf{Инженерия:} Python, FastAPI, asyncio, PostgreSQL / SQL, Redis, Docker, CI/CD, Linux, Git, pytest
  \item \textbf{Математика и алгоритмы:} теория вероятностей, матстатистика, линейная алгебра, методы оптимизации;
        C++, алгоритмы и структуры данных, оценка сложности
\end{itemize}

% ===== Experience =====
\section{Опыт работы}

\role{ReachFlow --- ML/NLP \& Backend Engineer}{Янв 2026 --- наст. время}
\textit{B2B SaaS: ИИ-ассистент, который ведёт продающие переписки в Telegram \textperiodcentered{} \href{https://reachflow.tech}{reachflow.tech}}
\begin{itemize}
  \item Построил \textbf{LLM-конвейер диалога} --- 8 вызовов модели на входящее сообщение: гейт батча,
        классификатор состояния лида (16 типов событий с якорем на сообщение и confidence), извлечение и взвешивание
        фактов (18 типизированных ключей), стратег\,$\rightarrow$\,гуманизатор\,$\rightarrow$\,ревьюер, планировщик отправки.
  \item \textbf{Снизил стоимость хода}: слил пять классификаторных вызовов в два --- одна модель отдаёт несколько
        секций одним объектом; 10 вызовов на сообщение стало 8, входных символов на стадии классификации на
        14--43\% меньше, при этом каждая секция сохранила собственный режим отказа.
  \item Сделал \textbf{strict structured output}: реестр из 25 LLM-задач и 22 JSON-схем на pydantic, нормализация
        схем под strict-режим и валидация ответа на своей стороне (нарушение схемы --- retryable); исполнитель с
        ретраями, каскадом фолбэков между вендорами и трейсом каждой попытки (модель, latency, retry, стоимость).
  \item Собрал \textbf{ИИ-агента с инструментами}: реестр \texttt{ToolSpec} --- единственный источник правды для блока
        инструментов в промпте, гейта привилегий по риску, границы доверия и MCP-эндпоинта; оркестратор с бюджетом вызовов.
  \item Разобрал провал \textbf{семантического гейта}: косинусная близость на эмбеддингах (text-embedding-3-small,
        pgvector / HNSW) отбраковывала 55 кандидатов из 63, потому что мерила формат и длину текста, а не
        релевантность --- перевёл гейт из блокирующего в наблюдательный.
\end{itemize}
\vspace{3pt}

\role{Freelance Python Developer}{Июль 2024 --- наст. время}
\begin{itemize}
  \item Telegram-бот-консультант для дистрибьютора пластиковых окон: \textbf{RAG} на PostgreSQL с векторным поиском
        по базе знаний, интеграция с CRM Bitrix24; по оценке заказчика закрыл ${\sim}80\%$ типовых консультаций.
\end{itemize}
\vspace{3pt}

\role{Near AI --- разработчик решений на C++}{Июнь 2023 --- Июль 2024}
\begin{itemize}
  \item Алгоритмические решения на C++ как обучающие данные для AI-платформы: оптимизация по времени и памяти,
        разбор асимптотики, code review в международной команде.
\end{itemize}

% ===== ML projects =====
\section{ML-проекты и соревнования}

\role{AIDAO 2025 --- Международная олимпиада по ИИ и анализу данных (ФКН НИУ ВШЭ, Яндекс Образование)}{Ноя 2025}
\textit{Финал: топ-30 команд из 248 из 14 стран, 32 часа очно. Задача команды беспилотных технологий Яндекса ---
occupancy-карта в bird's-eye view по камерам автомобиля, метрика IoU. PyTorch, решение с нуля.}
\begin{itemize}
  \item Реализовал \textbf{Lift--Splat--Shoot}: общий для камер ResNet-50 с FPN-слиянием уровней /8 и /16,
        depth-head на 64 бина глубины с обучаемой температурой softmax, дифференцируемый splat признаков в
        BEV-сетку 188$\times$126, фьюжн камер max\,+\,mean с SE-блоком, декодер на dilated residual-блоках.
  \item Основной прирост дала \textbf{покадровая калибровка} вместо фиксированной: геометрия проекции пересчитывается
        на каждом кадре из \textit{intrinsics} и \textit{car\_to\_cam}, с переносом сетки пикселей в координаты
        исходного изображения.
  \item Лосс --- BCE с \texttt{pos\_weight} по балансу классов $+\ 0.3\times$ дифференцируемый surrogate mIoU:
        оптимизировал ровно ту метрику, по которой считался скор. AdamW с раздельными lr для backbone и голов,
        cosine annealing, AMP, калибровка порога свипом по валидации.
  \item Итог --- \textbf{IoU 0.5606} на скрытом тесте (у команды-победителя олимпиады 0.564); инференс 2000 кадров
        за 1{,}7 мин на одной GPU.
\end{itemize}

% ===== Education =====
\section{Образование и курсы}
\textbf{Московский политехнический университет} --- прикладная математика и информатика, выпуск 2029
\textperiodcentered{} \textbf{НИУ ВШЭ, ФКН} --- вольнослушатель курсов ПМИ
\textperiodcentered{} \textbf{Deep Learning School} (ФПМИ МФТИ), 2025
\textperiodcentered{} \textbf{ОЦ <<Сириус>>} --- Январская научная школа по математике и программированию, 92 ак. ч., 2025
\textperiodcentered{} \textbf{Т-Поколение} (Т-Банк) --- <<Алгоритмы>>, параллель B, 2025

\end{document}
